\documentclass[12pt]{article}
\usepackage[utf8]{inputenc}
\usepackage{amsmath, amssymb, amsfonts, bm}
\usepackage{geometry}
\geometry{margin=1in}
\usepackage{enumitem}
\setlist[enumerate]{leftmargin=*}

\title{Solutions -- Quiz 1\\ \vspace{8pt} \Large ECON 308 -- International Economic Relations}
\author{Instructor: Muhebullah Karimzada}
\date{September 18, 2025}
\usepackage{comment}


\begin{document}
\maketitle


\begin{enumerate}



\item \textbf{Data recap.} Unit labor requirements:
\[
\begin{array}{c|cc}
 & \text{Wine} & \text{Cheese}\\\hline
\text{Home} & a_{LW}=2 & a_{LC}=1\\
\text{Foreign} & a_{LW}^{*}=6 & a_{LC}^{*}=2
\end{array}
\qquad \text{Labor: }L=L^{*}=120.
\]

\begin{enumerate}[label=(\alph*)]
\item \textbf{PPF (maximum production possibilities).}
\begin{itemize}
  \item \emph{Home:} With $L=120$, the maximum Wine is $W^{\max}=120/2=60$, and the maximum Cheese is $C^{\max}=120/1=120.$ The PPF is
  \[
  2W + 1\cdot C = 120 \quad \Longleftrightarrow \quad C = 120 - 2W.
  \]
  \item \emph{Foreign:} With $L^{*}=120$, the maximum Wine is $W^{\max*}=120/6=20$, and the maximum Cheese is $C^{\max*}=120/2=60.$ The PPF is
  \[
  6W + 2\cdot C = 120 \quad \Longleftrightarrow \quad C = 60 - 3W.
  \]
\end{itemize}

\item \textbf{Opportunity cost of Wine (in units of Cheese).}
In the Ricardian model, the opportunity cost of Wine in terms of Cheese equals the ratio of unit labor requirements:
\[
\text{OC}_W^{\text{Home}} = \frac{a_{LW}}{a_{LC}} = \frac{2}{1} = 2 \text{ (Cheese per Wine)},\qquad
\text{OC}_W^{\text{Foreign}} = \frac{a_{LW}^*}{a_{LC}^*} = \frac{6}{2} = 3.
\]
Equivalently, the opportunity cost of Cheese in terms of Wine is the inverse:
\[
\text{OC}_C^{\text{Home}} = \frac{1}{2},\qquad \text{OC}_C^{\text{Foreign}} = \frac{1}{3}.
\]

\item \textbf{Comparative advantage.}
A country has comparative advantage in the good with the lower opportunity cost:
\[
\text{OC}_W^{\text{Home}}=2 < 3=\text{OC}_W^{\text{Foreign}} \Rightarrow \text{Home has C.A. in Wine.}
\]
\[
\text{OC}_C^{\text{Foreign}}=\tfrac{1}{3} < \tfrac{1}{2}=\text{OC}_C^{\text{Home}} \Rightarrow \text{Foreign has C.A. in Cheese.}
\]

\item \textbf{Trade pattern at world relative price $P_W/P_C=1.5$.}
In a one-factor Ricardian model, a country produces the good for which the world relative price exceeds its autarky unit-cost ratio (its opportunity cost). Home will produce/export Wine if
\[
\frac{P_W}{P_C} > \frac{a_{LW}}{a_{LC}} = 2.
\]
But $1.5 < 2$, so Home does \emph{not} produce Wine; it produces Cheese.

Foreign will produce/export Wine if
\[
\frac{P_W}{P_C} > \frac{a_{LW}^*}{a_{LC}^*} = 3.
\]
But $1.5 < 3$, so Foreign also produces Cheese.

\textit{Conclusion:} At $P_W/P_C=1.5$, Wine is relatively cheap compared to both countries' opportunity costs, so \emph{both} countries supply Cheese and import Wine (from the rest of the world). If we restrict the world to only these two countries, such a price cannot be an equilibrium; an equilibrium relative price must lie between $2$ and $3$ for mutually consistent specialization (Home in Wine, Foreign in Cheese).
\end{enumerate}

\bigskip
\hrule
\bigskip


\item \textbf{Gravity equation:} 
\[
T_{ij} = A \cdot \frac{Y_i Y_j}{D_{ij}}, \quad A=1.
\]
\textbf{Data:} $Y_A=2000,\; Y_B=1500,\; Y_C=500$ (all in billions), $D_{AB}=1000$ km, $D_{AC}=5000$ km.

\begin{enumerate}[label=(\alph*)]
\item \textbf{Predicted trade $T_{AB}$ and $T_{AC}$.}
\[
T_{AB}=\frac{(2000)(1500)}{1000}=\frac{3{,}000{,}000}{1000}=3000.
\]
\[
T_{AC}=\frac{(2000)(500)}{5000}=\frac{1{,}000{,}000}{5000}=200.
\]
\[
\boxed{T_{AB}=3000,\quad T_{AC}=200 \quad (\text{in proportional units}).}
\]

\item \textbf{Trade ratio $T_{AB}/T_{AC}$.}
\[
\frac{T_{AB}}{T_{AC}}=\frac{3000}{200}=15.
\]
\[
\boxed{T_{AB} \text{ is 15 times } T_{AC}.}
\]

\item \textbf{FTA halves effective distance for A--C: } $D_{AC}'=2500$ km. Recompute:
\[
T_{AC}'=\frac{(2000)(500)}{2500}=\frac{1{,}000{,}000}{2500}=400.
\]
\[
\boxed{T_{AC}'=400.}
\]
\textit{Effect:} A halving of effective distance doubles predicted trade (from $200$ to $400$). The A--B vs.\ A--C ratio falls from $15$ to $3000/400=7.5$, indicating a substantial relative increase in A--C trade intensity.
\end{enumerate}
\end{enumerate}
\end{document}
